\part{Principii de programare}

Programarea este o activitate profund socială. De cînd există software, oamenii au schimbat cod între ei și au învățat unii de la alții cum să scrie cod. Programele pot trece prin mîna mai multor programatori concomitent și succesiv, uneori continuînd să existe la zeci de ani după ce dezvoltatorii originali au părăsit echipa.

Prin contrast, programărea competitivă este o activitate mai singuratică. De cele mai multe ori programele sînt dezvoltate individual sau în echipe mici, sînt scurte și sînt efemere. Rostul lor primar dispare odată cu încheierea concursului. Acest stil duce la obiceiuri de programare foarte diferite și adesea nocive.

Este fantastic că jurizarea programelor este automatizată, rapidă și ferită de slăbiciunea umană. Aceasta ne diferențiază de concursurile altor materii. Dar, tocmai de aceea, profesorii tind să se declare mulțumiți cînd elevii iau 100 de puncte sau \textit{Accepted} pe vreun site, fără să le citească niciodată codul-sursă.

În plus, una dintre părțile cele mai frumoase ale algoritmicii este că algoritmii sînt corecți în mod demonstrabil. Paradoxal, tocmai asta duce la cod mai urît, căci nici profesorii, nici elevii nu prea pun preț pe detaliile de implementare și pe lizibilitate.

Nu în ultimul rînd, la nivel de liceu există o ruptură între mediul academic și mediul industrial (la facultate legătura este cu mult mai strînsă). Majoritatea profesorilor de informatică de liceu nu au citit și nu au scris niciodată cod practic și nu simt nevoia să codeze clar.

Sper ca în această parte să vă conving că stilul de codare, precum și adoptarea unor bune practici de inginerie software considerate normale în industrie, vă pot ajuta și în programarea competitivă. Nu este obligatoriu să începeți cu această parte, dar sper să reveniți curînd la ea. Dacă ați ales să dedicați niște zeci sau sute de ore acestei cărți, înseamnă că vă doriți să creșteți ca programatori, iar atunci lăsați-mă și pe mine să decid cum se va întîmpla asta, nu alegeți doar ce credeți voi că este suficient. \emoji{slightly-smiling-face}

\import{./principles}{spaghetti-code.tex}
